\documentclass[a4paper,11pt]{article}

\usepackage{revy}
\usepackage[utf8]{inputenc}
\usepackage[T1]{fontenc}
\usepackage[danish]{babel}
\usepackage{amsmath,amssymb}

\revyname{MatRevy}
\revyyear{2026}
\version{2026-10-09}
\eta{1 minutter}

\title{Ofringen}
\author{Lea' 23 \& Carl '21}

\begin{document}
\maketitle

\begin{roles}
\role{Sx}[Statister] Kridt i munkekåber
\end{roles}

\begin{sketch}
\scene{Lys op. På scene er der 4-7 kridt i munkekåber igang med et opfrelsesritual. Der ligger et kridt i midten bundet fast, de andre står i en halvcirkel rundt om med lys. Siger måske nogle lyde oooo-meee-gaaa x10, så ryster scenen. Der kommer en kæmpe hånd (rekvisit) ud bag bagtæppet. Kridtet i midten begynder at lave bange lyde (lea demonstrerer). Munkekridtene begynder at bære kridtet hen mod hånden, som måske begynder at skrige eller kæmpe imod. Kridtet når hen til hånden, som bærer kridtet ud.}

\scene{Video: På AV ser man en kort video, af en person der tager et kridt op af lommen, tager det op mod tavlen, kridtet knækker...}
\end{sketch}
\end{document}
